\documentclass[a4paper,11pt]{extarticle}
\usepackage{fontspec}\setmainfont{DejaVu Serif}[Scale=0.9]
\usepackage{polyglossia}\setmainlanguage{russian}
\usepackage{amsmath,amssymb,array,booktabs,geometry,enumitem,xcolor}
\usepackage[most]{tcolorbox}
\geometry{left=1.7cm,right=1.7cm,top=1.5cm,bottom=1.5cm}
\pagestyle{plain}\setlength{\parindent}{0pt}\setlength{\parskip}{4pt}
\newcommand{\kr}{\text{кр}}\newcommand{\nab}{\text{набл}}
\setlist{nosep,leftmargin=1.3em}
\newtcolorbox{idea}{colback=green!6,colframe=green!45!black,boxrule=0.6pt,arc=2pt,left=5pt,right=5pt,top=3pt,bottom=3pt,title={Главная мысль},fonttitle=\bfseries\small}
\newtcolorbox{warn}{colback=orange!8,colframe=orange!70!black,boxrule=0.6pt,arc=2pt,left=5pt,right=5pt,top=3pt,bottom=3pt,title={Нюанс},fonttitle=\bfseries\small}
\newtcolorbox{task}[1]{colback=blue!4,colframe=blue!50!black,boxrule=0.6pt,arc=2pt,left=5pt,right=5pt,top=3pt,bottom=3pt,title={#1},fonttitle=\bfseries\small,breakable}
\newtcolorbox{tmpl}{colback=white,colframe=black,boxrule=0.8pt,arc=0pt,left=6pt,right=6pt,top=4pt,bottom=4pt}
\begin{document}
\begin{center}{\LARGE\bfseries Матстат с нуля: минимум теории для самостоятельной}\\[3pt]
{\small Пишите всё руками в тетрадь: формулы, таблицы, шаги. Задачи --- в тетради, ответы в конце файла.}\end{center}

\section*{0. Как работать с этим файлом}
\begin{enumerate}
\item Читайте по разделам и \textbf{переписывайте} рамки «Главная мысль» и «Нюанс», формулы и шаблоны.
\item После каждого раздела решайте мини-задачи (синие рамки). Подсказки в них --- это буквально какую формулу взять.
\item Ответы с разбором --- в разделе 9. Сначала решить, потом смотреть.
\item Все критические точки для задач берутся из таблиц в файле \emph{Шпаргалка\_критерии.pdf} (раздел 5 там).
\end{enumerate}

\section*{1. Словарь обозначений (выучить, это язык всех задач)}
\renewcommand{\arraystretch}{1.25}
\noindent\begin{tabular}{|>{\centering\arraybackslash}p{2.4cm}|p{8.2cm}|p{5.2cm}|}\hline
\textbf{Символ} & \textbf{Что это} & \textbf{Откуда берётся}\\\hline
$n$ & объём выборки (сколько измерили) & дано\\\hline
$\bar x$ & выборочное среднее & считаем по данным\\\hline
$s^2$, $s$ & \emph{исправленная} дисперсия и $s=\sqrt{s^2}$ (разброс в выборке) & считаем по данным или дано\\\hline
$\mu_0$ (или $a_0$, $m_0$) & число, которое \emph{утверждает} $H_0$ про среднее & дано в условии\\\hline
$\sigma^2=D(X)$ & \emph{истинная} дисперсия генеральной совокупности & дано, если «известна»\\\hline
$\sigma_0^2$ & число, которое утверждает $H_0$ про дисперсию & дано в условии\\\hline
$H_0$, $H_1$ & основная гипотеза («ничего не изменилось») и конкурирующая & формулируем сами по условию\\\hline
$\alpha$ & уровень значимости (допустимый риск ложной тревоги) & дано (0,05; 0,01; 0,1)\\\hline
$k$ & число степеней свободы (для таблиц) & считаем по формуле для критерия\\\hline
$\Phi(z)$ & функция Лапласа (нужна для таблицы нормального закона) & таблица\\\hline
\end{tabular}

\section*{2. Выборка, среднее и разброс}
\textbf{Генеральная совокупность} --- «всё» (вся партия деталей, все пакеты сока). \textbf{Выборка} --- то, что реально измерили ($n$ штук). По выборке мы судим о всей совокупности. Статистика --- это про то, как не ошибиться в таких выводах.

\textbf{Среднее:} $\displaystyle\bar x=\frac{x_1+x_2+\dots+x_n}{n}=\frac{\sum x_i}{n}$.

\textbf{Разброс.} Среднее мало что говорит: у выборок $4,6,8,10,12$ и $8,8,8,8,8$ среднее одно и то же (8), но разброс разный. Меру разброса строят так: посмотреть, на сколько каждое число отличается от среднего ($x_i-\bar x$), возвести в квадрат (чтобы минусы не гасили плюсы), сложить и разделить:
$$s^2=\frac{\sum(x_i-\bar x)^2}{n-1},\qquad s=\sqrt{s^2}.$$

\begin{warn}
\textbf{Почему $n-1$, а не $n$.} Выборка в среднем \emph{занижает} настоящий разброс, поэтому делят на $n-1$. Такая дисперсия называется \textbf{исправленной}. Во всех формулах критериев нужна именно она. Если в условии дана «выборочная» $D_{\text{в}}$ (делили на $n$), переходим: $s^2=\dfrac{n}{n-1}D_{\text{в}}$.
\end{warn}
\begin{warn}
\textbf{Две «дисперсии» --- не путать.} $s^2$ \emph{считаем по выборке} (оценка). $\sigma^2=D(X)$ --- \emph{истинная} для всей совокупности, её дают в условии со словами «известна». Если дано СКО $\sigma$, а нужна дисперсия, возводим в квадрат; если дана дисперсия, а нужно СКО, берём корень.
\end{warn}
\textbf{Удобная таблица для расчёта} (рисуем всегда): столбцы $x_i$ \,|\, $x_i-\bar x$ \,|\, $(x_i-\bar x)^2$. Проверка: сумма второго столбца равна 0.

\begin{task}{Мини-задача A1}
Выборка: $4,\,6,\,8,\,10,\,12$. Найдите $\bar x$, составьте таблицу отклонений, найдите $s^2$ и $s$.\\
\emph{Наводка:} $\bar x=\sum x_i/n$; затем $x_i-\bar x$ и квадраты; $s^2=\sum(\dots)^2/(n-1)$.
\end{task}
\begin{task}{Мини-задача A2}
(а) Дано $D(X)=2{,}25$. Чему равно СКО $\sigma$? \quad (б) Дана выборочная дисперсия $D_{\text{в}}=3{,}6$ при $n=10$. Найдите исправленную $s^2$.
\end{task}

\section*{3. Нормальное распределение и функция Лапласа}
Большинство величин в задачах (размеры деталей, вес, ошибки измерений) подчиняются \textbf{нормальному закону} («колокол»): чаще всего значения близки к среднему $\mu$, и чем дальше, тем реже. Разброс колокола задаёт $\sigma$. Запомните картинку:
\begin{center}
\begin{tabular}{|c|c|}\hline
в пределах $\mu\pm1\sigma$ & около 68\% значений\\\hline
в пределах $\mu\pm1{,}96\sigma$ & \textbf{95\%} значений\\\hline
в пределах $\mu\pm2{,}58\sigma$ & \textbf{99\%} значений\\\hline
\end{tabular}
\end{center}
Отсюда \textbf{откуда берутся числа 1,96 и 2,58}: если $H_0$ верна, то «обычное» значение попадает внутрь $\pm1{,}96$ с вероятностью 95\%, а редкие 5\% лежат снаружи. Именно эти 5\% (то есть $\alpha=0{,}05$) мы и считаем «подозрительными».

\textbf{Стандартизация.} Любое нормальное значение можно перевести в «стандартную линейку»:
$$z=\frac{x-\mu}{\sigma}\quad(\text{на сколько сигм значение отошло от среднего}).$$
Все критерии в этом курсе --- вариации этой идеи.

\textbf{Функция Лапласа} $\Phi(z)$ --- площадь под колоколом от 0 до $z$ (то есть доля значений между центром и $z$). Работаем с таблицей в обе стороны: \quad $\Phi(1{,}96)=0{,}475$,\quad $\Phi(2{,}58)=0{,}495$,\quad $\Phi(1{,}645)=0{,}45$. Мы \textbf{не вычисляем} $\Phi$, а \textbf{ищем по таблице}: нашли число $0{,}475$ внутри таблицы $\Rightarrow$ считали с краёв $z=1{,}96$.

\begin{task}{Мини-задача A3}
Найдите $z$ по таблице Лапласа (в шпаргалке п.\,5.1 есть нужные значения): (а) $\Phi(z)=0{,}475$; (б) $\Phi(z)=0{,}495$; (в) $\Phi(z)=0{,}45$.
\end{task}
\begin{task}{Мини-задача A4}
Размер детали $X\sim N(\mu=50,\ \sigma=4)$. Деталь имеет размер $58$. Найдите $z=\dfrac{x-\mu}{\sigma}$. Это «обычное» значение или подозрительное (сравните с 1,96)?
\end{task}

\section*{4. Идея проверки гипотезы (без формул)}
\textbf{История.} Автомат должен наливать 1000 мл. Взяли 25 пакетов, среднее вышло 996. Это сбой автомата или просто случайность? В выборке и должно быть небольшое отклонение. Вопрос --- насколько оно «слишком большое».

\textbf{Как это решают:}
\begin{enumerate}
\item Выдвигаем \textbf{$H_0$ --- «всё в порядке, отклонение случайно»}: $M(X)=1000$. Выдвигаем $H_1$ --- что именно мы подозреваем.
\item Считаем \textbf{наблюдаемое значение} $\;=\;\dfrac{\text{что получили}-\text{что утверждает }H_0}{\text{типичная случайная ошибка}}$. Это число «в сигмах»: насколько странный у нас результат.
\item Если $H_0$ верна, такие числа должны быть небольшими (в границах $\pm1{,}96$ и т.\,п.). Границу задаёт \textbf{критическая точка}, она зависит от $\alpha$ (насколько мы осторожны).
\item Наблюдаемое значение \textbf{за границей} (в критической области) $\Rightarrow$ слишком странно для случайности $\Rightarrow$ \textbf{отвергаем $H_0$}. Внутри $\Rightarrow$ \textbf{нет оснований отвергнуть} $H_0$.
\end{enumerate}
\begin{idea}
Все критерии устроены одинаково: \textbf{считаем число, сравниваем с границей из таблицы}. Меняются только формула числа и таблица.
\end{idea}

\textbf{$\alpha$ --- уровень значимости} = вероятность \emph{ложной тревоги} (отвергнуть $H_0$, когда она на самом деле верна). $\alpha=0{,}01$ строже (труднее отвергнуть), чем $\alpha=0{,}05$. Чем меньше $\alpha$, тем дальше граница.

\textbf{Вид $H_1$ определяет, где лежит критическая область:}
\begin{center}
\begin{tabular}{|c|c|p{7.4cm}|}\hline
$H_1$ & область & когда уместно\\\hline
$\ne$ & двусторонняя (обе стороны) & «изменилось», «отличается» (в любую сторону); на каждый хвост $\alpha/2$\\\hline
$>$ & правосторонняя & «увеличилось», «больше», «превышает»\\\hline
$<$ & левосторонняя & «уменьшилось», «меньше», «недолив»\\\hline
\end{tabular}
\end{center}
\begin{warn}
\textbf{Отвергнуть} ≠ \textbf{доказать противоположное}. Если наблюдаемое значение не попало в критическую область, пишем строго: «нет оснований отвергнуть $H_0$». Фразу «$H_0$ верна» писать нельзя.
\end{warn}

\begin{task}{Мини-задача B1}
Напишите $H_0$ и $H_1$: (а) «номинал 25 мм; проверить, что средний размер равен номиналу (отличается ли он)»; (б) «норма 1000 мл; есть подозрение, что автомат \emph{недоливает}»; (в) «средний брак был 3\%; проверяем, \emph{вырос} ли он после обновления».
\end{task}
\begin{task}{Мини-задача B2}
По $\alpha=0{,}05$ найдите $z_\kr$ для критерия $Z$: (а) $H_1$ с $\ne$ (формула $\Phi(z_\kr)=\frac{1-\alpha}{2}$); (б) $H_1$ с $>$ (формула $\Phi(z_\kr)=\frac{1-2\alpha}{2}$).
\end{task}
\begin{task}{Мини-задача B3}
Вывод в критерии $Z$. (а) $Z_\nab=-2{,}3$, $H_1$: $\ne$, $\alpha=0{,}05$ ($z_\kr=1{,}96$). Отвергаем $H_0$? \quad (б) $Z_\nab=-2{,}3$, $H_1$: $<$, $\alpha=0{,}01$ ($z_\kr=2{,}33$, область $Z_\nab<-z_\kr$). Отвергаем $H_0$?
\end{task}

\section*{5. Шаблон решения (переписать на отдельный лист --- это ваш «бланк»)}
\begin{tmpl}
\textbf{Дано:} $n=\dots$, $\bar x=\dots$, (что известно про дисперсию: $\sigma^2$ дана / $s^2$ дана), $\alpha=\dots$\\
\textbf{Гипотезы:} $H_0:\ \dots$, \quad $H_1:\ \dots$ (знак определяет область)\\
\textbf{Шаг 1. Критерий} (по таблице выбора из раздела 6) и почему.\\
\textbf{Шаг 2. Степени свободы} $k=\dots$ (если нужны; для $Z$ их нет).\\
\textbf{Шаг 3. Критическая точка} из таблицы (для $Z$ --- через $\Phi$; для $t,\chi^2,F$ --- по $\alpha$ и $k$).\\
\textbf{Шаг 4. Наблюдаемое значение:} формула $\to$ подстановка $\to$ число.\\
\textbf{Шаг 5. Сравнение и вывод:} сравнить; «$H_0$ отвергается» или «нет оснований отвергнуть $H_0$»; вывод словами про задачу.
\end{tmpl}

\section*{6. Откуда формулы и как выбрать критерий}
\subsection*{6.1 Один «секрет» формул}
Все формулы наблюдаемого значения --- это \textbf{(получили $-$ утверждает $H_0$) / типичная случайная ошибка}:
\begin{center}
\renewcommand{\arraystretch}{1.5}
\begin{tabular}{|l|c|p{5.6cm}|}\hline
критерий & формула & что в знаменателе\\\hline
$Z$ (1 выборка) & $\dfrac{(\bar x-\mu_0)\sqrt n}{\sigma}=\dfrac{\bar x-\mu_0}{\sigma/\sqrt n}$ & типичная ошибка среднего: $\sigma/\sqrt n$. $\sqrt n$ --- потому что чем больше выборка, тем точнее среднее\\\hline
$t$ (1 выборка) & $\dfrac{(\bar x-\mu_0)\sqrt n}{s}$ & то же, но $\sigma$ неизвестна, подставили $s$\\\hline
$Z$ (2 выборки) & $\dfrac{\bar x-\bar y}{\sqrt{D(X)/n_1+D(Y)/n_2}}$ & ошибки двух средних складываются (дисперсии складываются)\\\hline
$\chi^2$ & $\dfrac{(n-1)s^2}{\sigma_0^2}$ & отношение «что получили» к «что утверждает $H_0$» (в «разах»)\\\hline
$F$ & $\dfrac{s^2_{\text{больш}}}{s^2_{\text{меньш}}}$ & сравниваем дисперсии отношением: если равны, $F\approx1$\\\hline
\end{tabular}
\end{center}
\textbf{Быстрая проверка здравым смыслом:} $Z$ и $T$ около $0$ $\Rightarrow$ всё нормально; $F$ около $1$ $\Rightarrow$ дисперсии похожи; $\chi^2$ около $k$ $\Rightarrow$ нормально. Если $Z$ вышло $50$ или $F$ вышло $0{,}3$, вы ошиблись в вычислении (в $F$ большая дисперсия всегда сверху!).

\subsection*{6.2 Зачем столько разных таблиц}
\begin{itemize}
\item $\mathbf{Z}$ --- нормальный колокол. Применяется, когда $\sigma$ \textbf{известна}. Таблица: Лаплас.
\item $\mathbf{t}$ (Стьюдент) --- когда вместо $\sigma$ приходится подставлять $s$ (мы её оценили по малой выборке, добавилась неопределённость). «Колокол» становится чуть шире, поэтому критические точки больше, чем для $Z$ (при $n=9$: 2,306 вместо 1,96). При больших $n$ почти совпадает с $Z$. Нужны $k$ и $\alpha$.
\item $\boldsymbol{\chi^2}$ --- для \emph{дисперсии} (там складываются квадраты, а квадрат всегда $\ge0$, колокола нет: распределение «с горбом слева и хвостом справа»). Нужны $k=n-1$ и $\alpha$.
\item $\mathbf{F}$ --- для \emph{отношения двух дисперсий}. Нужны два числа степеней свободы $k_1,k_2$.
\item \textbf{Бартлетт} --- как $F$, но когда дисперсий больше двух. Таблица $\chi^2$, $k=l-1$.
\end{itemize}

\subsection*{6.3 Как выбрать критерий (три вопроса по порядку)}
\begin{tabular}{|p{6.4cm}|p{4.2cm}|p{5cm}|}\hline
\textbf{1. Что сравниваем?} & \textbf{2. Сколько выборок?} & \textbf{3. $\sigma$ известна?}\\\hline
\textbf{среднее} $M(X)=\dots$ & 1 или 2 & да $\to$ $Z$; нет $\to$ $t$\\\hline
\textbf{дисперсию} $D(X)=\dots$ & 1 (против числа) & $\chi^2$\\\hline
\textbf{дисперсию} $D(X)=D(Y)$ & 2 & Фишер ($F$)\\\hline
\textbf{дисперсию} (много) & 3 и больше & Бартлетт\\\hline
\end{tabular}
\\[3pt]
\textbf{Признаки в условии:} «дисперсия (или $\sigma$) \emph{известна}» $\to$ $Z$. Даны $s$, $s^2$ или таблица измерений, а про $\sigma$ ничего $\to$ $t$. «Исправленная дисперсия $s^2$» и $n$ при проверке дисперсии $\to$ $\chi^2$.

\section*{7. Нюансы, на которых теряют баллы}
\begin{enumerate}
\item \textbf{Дисперсия или СКО.} Под корнем в $Z$ --- дисперсия $D$; если дано $\sigma$, возведите в квадрат.
\item \textbf{$s^2$ исправленная} ($n-1$). Из $D_{\text{в}}$: $s^2=\frac{n}{n-1}D_{\text{в}}$.
\item \textbf{$\ne$ $\Rightarrow$ $\alpha/2$ или $\Phi=\frac{1-\alpha}{2}$.} \textbf{$>$ или $<$ $\Rightarrow$ $\Phi=\frac{1-2\alpha}{2}$.} Для $t$ в двусторонней таблице берём $\alpha$ при $\ne$ и $2\alpha$ при $>$ или $<$.
\item \textbf{Левосторонняя область ($H_1$: $<$)}: сравниваем $Z_\nab<-z_\kr$ (минус!).
\item \textbf{Фишер:} большая дисперсия сверху, $k_1$ --- у неё. Для $\ne$ берём $\alpha/2$.
\item \textbf{$\chi^2$:} для $>$ берём $\chi^2(\alpha;k)$, для $<$ --- $\chi^2(1-\alpha;k)$; для $\ne$ две границы, отвергаем \emph{вне} отрезка.
\item \textbf{Вывод:} «$H_0$ отвергается» или «нет оснований отвергнуть $H_0$». Никогда «доказали $H_0$».
\item \textbf{Округление:} 3--4 знака в промежуточных вычислениях, итог до 2 знаков.
\end{enumerate}

\newpage
\section*{8. Задачи по шаблону (решайте в тетради строго по бланку из раздела 5)}
\begin{task}{C1 \quad $Z$, одна выборка}
Дисперсия $D(X)=16$ известна, $n=64$, $\bar x=52{,}5$. При $\alpha=0{,}05$ проверьте $H_0$: $M(X)=50$ при $H_1$: $M(X)\ne50$.\\
\emph{Наводка:} $\sigma=\sqrt{D}$; $z_\kr$ из $\Phi=\frac{1-\alpha}{2}$; $Z_\nab=\frac{(\bar x-\mu_0)\sqrt n}{\sigma}$.
\end{task}
\begin{task}{C2 \quad $t$, одна выборка}
Из нормальной совокупности $n=9$, $\bar x=47$, $s=6$ ($\sigma$ неизвестна). $\alpha=0{,}05$. $H_0$: $M(X)=50$, $H_1$: $M(X)\ne50$.\\
\emph{Наводка:} $k=n-1$; $t_\kr$ из двусторонней таблицы Стьюдента; $T_\nab=\frac{(\bar x-\mu_0)\sqrt n}{s}$; сравнивать $|T_\nab|$.
\end{task}
\begin{task}{C3 \quad $\chi^2$, дисперсия}
$n=21$, исправленная $s^2=30$. $\alpha=0{,}05$. $H_0$: $D(X)=20$, $H_1$: $D(X)>20$.\\
\emph{Наводка:} $k=n-1$; правосторонняя область; $\chi^2_\nab=\frac{(n-1)s^2}{\sigma_0^2}$.
\end{task}
\begin{task}{C4 \quad Фишер}
Две выборки: $n_1=10,\ s_X^2=12$; $n_2=8,\ s_Y^2=4$. $\alpha=0{,}05$. $H_0$: $D(X)=D(Y)$, $H_1$: $D(X)>D(Y)$.\\
\emph{Наводка:} большая сверху; $k_1=n_{\text{больш}}-1$, $k_2=n_{\text{меньш}}-1$; таблица $F$ для $\alpha=0{,}05$.
\end{task}
\begin{task}{C5 \quad $Z$, две выборки}
$D(X)=9$, $D(Y)=16$ известны; $n_1=36$, $n_2=64$; $\bar x=20$, $\bar y=18{,}5$. $\alpha=0{,}05$. $H_0$: $M(X)=M(Y)$, $H_1$: $M(X)\ne M(Y)$.\\
\emph{Наводка:} $Z_\nab=\dfrac{\bar x-\bar y}{\sqrt{D(X)/n_1+D(Y)/n_2}}$; $z_\kr$ как в C1.
\end{task}
\begin{task}{C6 \quad $t$, две выборки (потренируйте формулу)}
$n_1=n_2=5$, $\bar x=10$, $\bar y=8$, $s_X^2=4$, $s_Y^2=2$ (дисперсии неизвестны, но равны). $\alpha=0{,}05$. $H_0$: $M(X)=M(Y)$, $H_1$: $\ne$.\\
\emph{Наводка:} $k=n_1+n_2-2$; считайте по частям: сначала $\sqrt{(n_1-1)s_X^2+(n_2-1)s_Y^2}$, затем $\sqrt{\frac{n_1n_2(n_1+n_2-2)}{n_1+n_2}}$, потом $T_\nab=\frac{\bar x-\bar y}{(\text{первое})}\cdot(\text{второе})$.
\end{task}
\begin{task}{C7 \quad Только выбор критерия (без вычислений)}
Какой критерий (и какую $k$) применить: (а) $n=20$, дана $s^2$, проверяем $D(X)=5$; (б) средние двух выборок, дисперсии известны; (в) одна выборка, $\sigma$ не дана, дан ряд измерений, проверяем среднее; (г) четыре станка, сравниваем их дисперсии; (д) две выборки, сначала нужно выяснить, одинаковы ли дисперсии.
\end{task}

\newpage
\section*{9. Ответы и разборы (после самостоятельного решения!)}
\small
\textbf{A1.} $\bar x=\frac{4+6+8+10+12}{5}=\frac{40}{5}=8$. Отклонения: $-4,-2,0,2,4$ (сумма 0 ✓.). Квадраты: $16,4,0,4,16$, сумма $40$. $s^2=\frac{40}{5-1}=10$, $s=\sqrt{10}\approx3{,}16$. (Если бы делили на $n$: $8$.)\\
\textbf{A2.} (а) $\sigma=\sqrt{2{,}25}=1{,}5$. (б) $s^2=\frac{10}{9}\cdot3{,}6=4{,}0$.\\
\textbf{A3.} (а) $1{,}96$; (б) $2{,}58$; (в) $1{,}645$.\\
\textbf{A4.} $z=\frac{58-50}{4}=2$. Это больше $1{,}96$: значение «подозрительное» (попадает в крайние 5\%).\\
\textbf{B1.} (а) $H_0$: $M(X)=25$, $H_1$: $M(X)\ne25$. (б) $H_0$: $M(X)=1000$, $H_1$: $M(X)<1000$. (в) $H_0$: $M(X)=3$, $H_1$: $M(X)>3$.\\
\textbf{B2.} (а) $\Phi=\frac{1-0{,}05}{2}=0{,}475\Rightarrow z_\kr=1{,}96$. (б) $\Phi=\frac{1-2\cdot0{,}05}{2}=0{,}45\Rightarrow z_\kr=1{,}645$.\\
\textbf{B3.} (а) $|{-2{,}3}|=2{,}3>1{,}96\Rightarrow$ отвергаем. (б) нужно $Z_\nab<-2{,}33$, а $-2{,}3>-2{,}33$ $\Rightarrow$ в критическую область не попали, \emph{нет оснований отвергнуть} (при строгом $\alpha=0{,}01$ граница дальше).

\medskip
\textbf{C1.} Дано: $D=16\Rightarrow\sigma=4$; $n=64$; $\bar x=52{,}5$; $\mu_0=50$. Критерий: $Z$ (дисперсия известна, 1 выборка). $\Phi(z_\kr)=\frac{1-0{,}05}{2}=0{,}475\Rightarrow z_\kr=1{,}96$. $Z_\nab=\frac{(52{,}5-50)\cdot\sqrt{64}}{4}=\frac{2{,}5\cdot8}{4}=5{,}0$. $5{,}0>1{,}96\Rightarrow$ $H_0$ отвергается: среднее отличается от 50.\\
\textbf{C2.} $t$ (дисперсия неизвестна). $k=8$; $t_{\text{двуст}}(0{,}05;8)=2{,}306$. $T_\nab=\frac{(47-50)\cdot\sqrt9}{6}=\frac{-3\cdot3}{6}=-1{,}5$. $|{-1{,}5}|<2{,}306\Rightarrow$ нет оснований отвергнуть $H_0$.\\
\textbf{C3.} $\chi^2$. $k=20$; $\chi^2_\kr(0{,}05;20)=31{,}41$. $\chi^2_\nab=\frac{20\cdot30}{20}=30$. $30<31{,}41\Rightarrow$ нет оснований отвергнуть $H_0$.\\
\textbf{C4.} $s^2_{\text{больш}}=12$ (первая выборка), $k_1=9$; $k_2=7$. $F_\kr(0{,}05;9,7)=3{,}68$. $F_\nab=\frac{12}{4}=3$. $3<3{,}68\Rightarrow$ нет оснований отвергнуть $H_0$.\\
\textbf{C5.} $Z$ (дисперсии известны, 2 выборки). $z_\kr=1{,}96$. $\sqrt{\frac9{36}+\frac{16}{64}}=\sqrt{0{,}25+0{,}25}=\sqrt{0{,}5}\approx0{,}707$. $Z_\nab=\frac{20-18{,}5}{0{,}707}\approx2{,}12>1{,}96\Rightarrow$ $H_0$ отвергается.\\
\textbf{C6.} $k=5+5-2=8$; $t_\kr(0{,}05;8)=2{,}306$. $\sqrt{4\cdot4+4\cdot2}=\sqrt{24}\approx4{,}899$. $\sqrt{\frac{5\cdot5\cdot8}{10}}=\sqrt{20}\approx4{,}472$. $T_\nab=\frac{2}{4{,}899}\cdot4{,}472\approx1{,}83$. $1{,}83<2{,}306\Rightarrow$ нет оснований отвергнуть $H_0$.\\
\textbf{C7.} (а) $\chi^2$, $k=19$. (б) $Z$ для двух выборок (без $k$). (в) $t$ для одной выборки, $k=n-1$. (г) Бартлетт, $k=l-1=3$. (д) Фишер, $k_1=n_{\text{б}}-1$, $k_2=n_{\text{м}}-1$ (если $H_0$ не отвергнута, дальше можно применять $t$ для двух выборок).
\end{document}
